\Chapter{34}

\Verse{1}
{
    \zA{话说袭人见贾母王夫人等去后，便走来宝玉身边坐下，含泪问他：“怎么就打 到这步田地？” 宝玉叹气说道：“不过为那些事，问他做什么!只是下半截疼的很，
        你瞧瞧，打坏了那里？” 袭人听说，便轻轻的伸手进去，将中衣脱下，略动一动， 宝玉便咬着牙叫嗳哟，袭人连忙停住手：如此三四次，才褪下来了。 袭人看时，只
        见腿上半段青紫，都有四指阔的僵痕高起来。 袭人咬着牙说道：“我的娘，怎么下 这般的狠手!你但凡听我一句话，也不到这个分儿。}
}
{
    \eA{When Hsi Jen saw dowager lady Chia, Madame Wang and the other members of the
        family take their leave, our narrative says, she entered the room. and,
        taking a seat next to Pao-yü, she asked him, while she did all she could to
        hide her tears: "How was it that he beat you to such extremes?" Pao-yü
        heaved a sigh. "It was simply," he replied, "about those trifles. But what's
        the use of your asking me about them?}
}
{
}

\Verse{2}
{
    \zA{幸而没动筋骨，倘或打出个残 疾来，可叫人怎么样呢？” 正说着，只听丫鬟们说：“宝姑娘来了。” 袭人听见，知道穿不及中衣，便拿了 一床夹纱被替宝玉盖了。
        只见宝钗手里托着一丸药走进来，向袭人说道：“晚上把 这药用酒研开，替他敷上，把那淤血的热毒散开，就好了。” 说毕，递与袭人。 又 问：“这会子可好些？”
        宝玉一面道谢，说：“好些了。” 又让坐。 宝钗见他睁开眼 说话，不像先时，心中也宽慰了些，便点头叹道：“早听人一句话，也不至有今日。}
}
{
    \eA{The lower part of my body is so very sore! Do look and see where I'm
        bruised!" At these words, Hsi Jen put out her hand, and inserting it gently
        under his clothes, she began to pull down the middle garments. She had but
        slightly moved them, however, when Pao-yü ground his teeth and groaned
        "ai-ya." Hsi Jen at once stayed her hand.}
}
{
}

\Verse{3}
{
    \zA{别说老太太、太太心疼，就是我们看着，心里也――”刚说了半句，又忙咽住，不 觉眼圈微红，双腮带赤，低头不语了。 宝玉听得这话如此亲切，大有深意，忽见他
        又咽住不往下说，红了脸低下头含着泪只管弄衣带，那一种软怯娇羞、轻怜痛惜之 情，竟难以言语形容，越觉心中感动，将疼痛早已丢在九霄云外去了。 想道：“我
        不过挨了几下打，他们一个个就有这些怜惜之态，令人可亲可敬。 假若我一时竟别 有大故，他们还不知何等悲感呢。}
}
{
    \eA{It was after three or four similar attempts that she, at length, succeeded
        in drawing them down. Then looking closely, Hsi Jen discovered that the
        upper part of his legs was all green and purple, one mass of scars four
        fingers wide, and covered with huge blisters. Hsi Jen gnashed her teeth. "My
        mother!" she ejaculated, "how is it that he struck you with such a ruthless
        hand!}
}
{
}

\Verse{4}
{
    \zA{既是他们这样，我便一时死了，得他们如此，一 生事业纵然尽付东流，也无足叹惜了。” 正想着，只听宝钗问袭人道：“怎么好好的动了气，就打起来了？” 袭人便把
        焙茗的话悄悄说了。 宝玉原来还不知贾环的话，见袭人说出，方才知道； 因又拉上 薛蟠，惟恐宝钗沉心，忙又止住袭人道：“薛大哥从来不是这样，你们别混猜度。”
        宝钗听说，便知宝玉是怕他多心，用话拦袭人。 因心中暗暗想道：“打得这个形象， 疼还顾不过来，还这样细心，怕得罪了人。}
}
{
    \eA{Had you minded the least bit of my advice to you, things wouldn't have come
        to such a pass! Luckily, no harm was done to any tendon or bone; for had you
        been crippled by the thrashing you got, what could we do?" In the middle of
        these remarks, she saw the servant-girls come, and they told her that Miss
        Pao-ch'ai had arrived.}
}
{
}

\Verse{5}
{
    \zA{你既这样用心，何不在外头大事上做工 夫，老爷也欢喜了，也不能吃这样亏。 你虽然怕我沉心所以拦袭人的话，难道我就
        不知我哥哥素日恣心纵欲、毫无防范的那种心性吗？ 当日为个秦种还闹的天翻地覆， 自然如今比先又加利害了。”
        想毕，因笑道：“你们也不必怨这个怨那个，据我想， 到底宝兄弟素日肯和那些人来往，老爷才生气。 就是我哥哥说话不防头，一时说出
        宝兄弟来，也不是有心挑唆：一则也是本来的实话，二则他原不理论这些防嫌小事。}
}
{
    \eA{Hearing this, Hsi Jen saw well enough that she had no time to put him on his
        middle garments, so forthwith snatching a double gauze coverlet, she threw
        it over Pao-yü. This done, she perceived Pao-ch'ai walk in, her hands laden
        with pills and medicines.}
}
{
}

\Verse{6}
{
    \zA{袭姑娘从小儿只见过宝兄弟这样细心的人，何曾见过我哥哥那天不怕地不怕、心里 有什么口里说什么的人呢？” 袭人因说出薛蟠来，见宝玉拦他的话，早已明白自己
        说造次了，恐宝钗没意思； 听宝钗如此说，更觉羞愧无言。 宝玉又听宝钗这一番话， 半是堂皇正大，半是体贴自己的私心，更觉比先心动神移。
        方欲说话时，只见宝钗 起身道：“明日再来看你，好生养着罢。 方才我拿了药来，交给袭人，晚上敷上管 就好了。” 说着便走出门去。}
}
{
    \eA{"At night," she said to Hsi Jen, "take these medicines and dissolve them in
        wine and then apply them on him, and, when the fiery virus from that
        stagnant blood has been dispelled, he'll be all right again." After these
        directions, she handed the medicines to Hsi Jen. "Is he feeling any better
        now?" she proceeded to inquired. "Thanks!" rejoined Pao-yü. "I'm feeling
        better," he at the same time went on to say;}
}
{
}

\Verse{7}
{
    \zA{袭人赶着送出院外，说：“姑娘倒费心了。 改日宝二爷 好了，亲自来谢。” 宝钗回头笑道：“这有什么的？ 只劝他好生养着，别胡思乱想就 好了。
        要想什么吃的玩的，悄悄的往我那里只管取去，不必惊动老太太、太太众人。 倘或吹到老爷耳朵里，虽然彼时不怎么样，将来对景，终是要吃亏的。” 说着去了。
        袭人抽身回来，心内着实感激宝钗。 进来见宝玉沉思默默，似睡非睡的模样， 因而退出房外栉沐。}
}
{
    \eA{after which, he pressed her to take a seat. Pao-ch'ai noticed that he could
        open his eyes wide, that he could speak and that he was not as bad as he had
        been, and she felt considerable inward relief. But nodding her head, she
        sighed. "If you had long ago listened to the least bit of the advice
        tendered to you by people things would not have reached this climax to-day,"
        she said.}
}
{
}

\Verse{8}
{
    \zA{宝玉默默的躺在床上，无奈臀上作痛，如针挑刀挖一般，更热 如火炙，略展转时，禁不住“嗳哟”之声。 那时天色将晚，因见袭人去了，却有两
        三个丫鬟伺候，此时并无呼唤之事，因说道：“你们且去梳洗，等我叫时再来。” 众 人听了，也都退出。
        这里宝玉昏昏沉沉，只见蒋玉函走进来了，诉说忠顺府拿他之事； 一时又见金 钏儿进来，哭说为他投井之情。 宝玉半梦半醒，刚要诉说前情，忽又觉有人推他，
        恍恍惚惚听得悲切之声。}
}
{
    \eA{"Not to speak of the pain experienced by our dear ancestor and aunt Wang,
        the sight of you in this state makes even us feel at heart…." Just as she
        had uttered half of the remark she meant to pass, she quickly suppressed the
        rest; and smitten by remorse for having spoken too hastily, she could not
        help getting red in the face and lowering her head.}
}
{
}

\Verse{9}
{
    \zA{宝玉从梦中惊醒，睁眼一看，不是别人，却是黛玉。 犹恐 是梦，忙又将身子欠起来，向脸上细细一认，只见他两个眼睛肿得桃儿一般，满面 泪光，不是黛玉却是那个？
        宝玉还欲看时，怎奈下半截疼痛难禁，支持不住，便“嗳 哟”一声仍旧倒下，叹了口气说道：“你又做什么来了？ 太阳才落，那地上还是怪热
        的，倘或又受了暑，怎么好呢？ 我虽然捱了打，却也不很觉疼痛。 这个样儿是装出 来哄他们，好在外头布散给老爷听。 其实是假的，你别信真了。”}
}
{
    \eA{Pao-yü was realising how affectionate, how friendly and how replete with
        deep meaning were the sentiments that dropped from her month, when, of a
        sudden, he saw her seal her lips and, flashing crimson, droop her head, and
        simply fumble with her girdle.}
}
{
}

\Verse{10}
{
    \zA{此时黛玉虽不是嚎啕大哭，然越是这等无声之泣，气噎喉堵，更觉利害。 听了 宝玉这些话，心中提起万句言词，要说时却不能说得半句。 半天，方抽抽噎噎的道：
        “你可都改了罢！” 宝玉听说，便长叹一声道：“你放心。 别说这样话。 我便为这些 人死了，也是情愿的。” 一句话未了，只见院外人说：“二奶奶来了。”
        黛玉便知是凤姐来了，连忙立起 身，说道：“我从后院子里去罢，回来再来。” 宝玉一把拉住道：“这又奇了，好好 的怎么怕起他来了？”}
}
{
    \eA{Yet so fascinating was she in those timid blushes, which completely baffle
        description, that his feelings were roused within him to such a degree, that
        all sense of pain flew at once beyond the empyrean. "I've only had to bear a
        few blows," he reflected, "and yet every one of them puts on those pitiful
        looks sufficient to evoke love and regard;}
}
{
}

\Verse{11}
{
    \zA{黛玉急得跺脚，悄悄的说道：“你瞧瞧我的眼睛!又该他们拿 咱们取笑儿了。” 宝玉听说，赶忙的放了手。 黛玉三步两步转过床后，刚出了后院，
        凤姐从前头已进来了。 问宝玉：“可好些了？ 想什么吃？ 叫人往我那里取去。” 接着薛 姨妈又来了。 一时贾母又打发了人来。
        至掌灯时分，宝玉只喝了两口汤，便昏昏沉沉的睡去。 接着周瑞媳妇、吴新登 媳妇、郑好时媳妇这几个有年纪长来往的，听见宝玉捱了打，也都进来。}
}
{
    \eA{so were, after all, any mishap or untimely end to unexpectedly befall me,
        who can tell how much more afflicted they won't be! And as they go on in
        this way, I shall have them, were I even to die in a moment, to feel so much
        for me; so there will indeed be no reason for regret, albeit the concerns of
        a whole lifetime will be thus flung entirely to the winds!"}
}
{
}

\Verse{12}
{
    \zA{袭人忙迎 出来，悄悄的笑道：“婶娘们略来迟了一步，二爷睡着了。” 说着，一面陪他们到那 边屋里坐着，倒茶给他们吃。
        那几个媳妇子都悄悄的坐了一回，向袭人说：“等二 爷醒了，你替我们说罢。” 袭人答应了，送他们出去。 刚要回来，只见王夫人使个
        老婆子来说：“太太叫一个跟二爷的人呢。” 袭人见说，想了一想，便回身悄悄的告 诉晴雯、麝月、秋纹等人说：“太太叫人，你们好生在屋里，我去了就来。”
        说毕， 同那老婆子一径出了园子，来至上房。}
}
{
    \eA{While indulging in these meditations, ha overheard Pao-ch'ai ask Hsi Jen:
        "How is it that he got angry, without rhyme or reason, and started beating
        him?" and Hsi Jen tell her, in reply, the version given to her by Pei Ming.
        Pao-yü had, in fact, no idea as yet of what had been said by Chia Huan, and,
        when he heard Hsi Jen's disclosures, he eventually got to know what it was;}
}
{
}

\Verse{13}
{
    \zA{王夫人正坐在凉榻上，摇着芭蕉扇子。 见他来了，说道：“你不管叫谁来也罢 了，又撂下他来了，谁伏侍他呢？” 袭人见说，连忙陪笑回道：“二爷才睡了，那
        四五个丫头，如今也好了，会伏侍了。 太太请放心。 恐怕太太有什么话吩咐，打发 他们来，一时听不明白倒耽误了事。”
        王夫人道：“也没什么话，白问问他这会子疼 的怎么样了？” 袭人道：“宝姑娘送来的药，我给二爷敷上了，比先好些了。 先疼
        的躺不住，这会子都睡沉了，可见好些。”}
}
{
    \eA{but as it also criminated Hsüeh P'an, he feared lest Pao-ch'ai might feel
        unhappy, so he lost no time in interrupting Hsi Jen. "Cousin Hsüeh," he
        interposed, "has never been like that; you people mustn't therefore give way
        to idle surmises!" These words were enough to make Pao-ch'ai see that Pao-yü
        had thought it expedient to say something to stop Hsi Jen's mouth,
        apprehending that her suspicions might get roused;}
}
{
}

\Verse{14}
{
    \zA{王夫人又问：“吃了什么没有？” 袭人道： “老太太给的一碗汤，喝了两口，只嚷干渴，要吃酸梅汤。 我想酸梅是个收敛东西，
        刚才捱打，又不许叫喊，自然急的热毒热血未免存在心里。 倘或吃下这个去激在心 里，再弄出病来，那可怎么样呢。 因此我劝了半天，才没吃。 只拿那糖腌的玫瑰卤
        子和了，吃了小半碗，嫌吃絮了，不香甜。” 王夫人道：“嗳哟，你何不早来和我说？ 前日倒有人送了几瓶子香露来。 原要给他一点子，我怕胡遭塌了，就没给。}
}
{
    \eA{and she consequently secretly mused within herself: "He has been beaten to
        such a pitch, and yet, heedless of his own pains and aches, he's still so
        careful not to hurt people's feelings. But since you can be so considerate,
        why don't you take a little more care in greater concerns outside, so that
        your father should feel a little happier, and that you also should not have
        to suffer such bitter ordeals!}
}
{
}

\Verse{15}
{
    \zA{既是他 嫌那玫瑰膏子吃絮了，把这个拿两瓶子去，一碗水里只用挑上一茶匙，就香的了不 得呢。” 说着，就唤彩云来：“把前日的那几瓶香露拿了来。”
        袭人道：“只拿两瓶来 罢，多也白遭塌。 等不够再来取也是一样。” 彩云听了，去了半日，果然拿了两瓶 来付与袭人。
        袭人看时，只见两个玻璃小瓶却有三寸大小，上面螺丝银盖，鹅黄笺 上写着“木樨清露”，那一个写着“玫瑰清露”。 袭人笑道：“好尊贵东西!这么个小
        瓶儿，能有多少？”}
}
{
    \eA{But notwithstanding that the dread of my feeling hurt has prompted you to
        interrupt Hsi Jen in what she had to tell me, is it likely that I am blind
        to the fact that my brother has ever followed his fancies, allowed his
        passions to run riot, and never done a thing to exercise any check over
        himself?}
}
{
}

\Verse{16}
{
    \zA{王夫人道：“那是进上的，你没看见鹅黄笺子？ 你好生替他收着， 别遭塌了。” 袭人答应着，方要走时，王夫人又叫：“站着，我想起一句话来问你。” 袭人忙
        又回来。 王夫人见房内无人，便问道：“我恍惚听见宝玉今日捱打，是环儿在老爷 跟前说了什么话，你可听见这个话没有？” 袭人道：“我倒没听见这个话，只听见
        说为二爷认得什么王府的戏子，人家来和老爷说了，为这个打的。” 王夫人摇头说 道：“也为这个。 只是还有别的原故呢。”}
}
{
    \eA{His temperament is such that he some time back created, all on account of
        that fellow Ch'in Chung, a rumpus that turned heaven and earth topsy-turvy;
        and, as a matter of course, he's now far worse than he was ever before!"
        "You people," she then observed aloud, at the close of these cogitations,
        "shouldn't bear this one or that one a grudge.}
}
{
}

\Verse{17}
{
    \zA{袭人道：“别的原故，实在不知道。” 又低 头迟疑了一会，说道：“今日大胆在太太跟前说句冒撞话，论理――”说了半截， 却又咽住。 王夫人道：“你只管说。”
        袭人道：“太太别生气，我才敢说。” 王夫人道： “你说就是了。” 袭人道：“论理宝二爷也得老爷教训教训才好呢!要老爷再不管，
        不知将来还要做出什么事来呢。” 王夫人听见了这话，便点头叹息，由不得赶着袭人叫了一声：“我的儿!你这话 说的很明白，和我的心里想的一样。}
}
{
    \eA{I can't help thinking that it's, after all, because of your usual readiness,
        cousin Pao-yü, to hobnob with that set that your father recently lost
        control over his temper. But assuming that my brother did speak in a
        careless manner and did casually allude to you cousin Pao-yü, it was with no
        design to instigate any one! In the first place, the remarks he made were
        really founded on actual facts;}
}
{
}

\Verse{18}
{
    \zA{其实，我何曾不知道宝玉该管？ 比如先时你珠 大爷在，我是怎么样管他，难道我如今倒不知管儿子了？ 只是有个原故：如今我想
        我已经五十岁的人了，通共剩了他一个，他又长的单弱，况且老太太宝贝似的，要 管紧了他，倘或再有个好歹儿，或是老太太气着，那时上下不安，倒不好，所以就
        纵坏了他了。 我时常掰着嘴儿说一阵，劝一阵，哭一阵。 彼时也好，过后来还是不 相干，到底吃了亏才罢!设若打坏了，将来我靠谁呢！”
        说着，由不得又滴下泪来。}
}
{
    \eA{and secondly, he's not one to ever trouble himself about such petty trifles
        as trying to guard against animosities. Ever since your youth up, Miss Hsi,
        you've simply had before your eyes a person so punctilious as cousin Pao-yü,
        but have you ever had any experience of one like that brother of mine, who
        neither fears the powers in heaven or in earth, and who readily blurts out
        all he thinks?"}
}
{
}

\Verse{19}
{
    \zA{袭人见王夫人这般悲感，自己也不觉伤了心，陪着落泪。 又道：“二爷是太太 养的，太太岂不心疼； 就是我们做下人的，伏侍一场，大家落个平安，也算造化了。
        要这样起来，连平安都不能了。 那一日那一时我不劝二爷？ 只是再劝不醒。 偏偏那 些人又肯亲近他，也怨不得他这样。 如今我们劝的倒不好了。
        今日太太提起这话来， 我还惦记着一件事，要来回太太，讨太太个主意。 只是我怕太太疑心，不但我的话 白说了，且连葬身之地都没有了！”}
}
{
    \eA{Hsi Jen, seeing Pao-yü interrupt her, at the bare mention of Hsüeh P'an,
        understood at once that she must have spoken recklessly and gave way to
        misgivings lest Pao-ch'ai might not have been placed in a false position,
        but when she heard the language used by Pao-ch'ai, she was filled with a
        keener sense of shame and could not utter a word.}
}
{
}

\Verse{20}
{
    \zA{王夫人听了这话内中有因，忙问道：“我的儿! 你只管说。 近来我因听见众人背前面后都夸你，我只说你不过在宝玉身上留心，或 是诸人跟前和气这些小意思。
        谁知你方才和我说的话，全是大道理，正合我的心事。 你有什么只管说什么，只别叫别人知道就是了。” 袭人道：“我也没什么别的说，我
        只想着讨太太一个示下，怎么变个法儿，以后竟还叫二爷搬出园外来住就好了。” 王夫人听了，吃一大惊，忙拉了袭人的手，问道：“宝玉难道和谁作怪了不成？”}
}
{
    \eA{Pao-yü too, after listening to the sentiments, which Pao-ch'ai expressed,
        felt, partly because they were so magnanimous and noble, and partly because
        they banished all misconception from his mind, his heart and soul throb with
        greater emotion then ever before. When, however, about to put in his word,
        he noticed Pao-ch'ai rise to her feet.}
}
{
}

\Verse{21}
{
    \zA{袭人连忙回道：“太太别多心，并没有这话，这不过是我的小见识：如今二爷也大 了，里头姑娘们也大了，况且林姑娘宝姑娘又是两姨姑表姐妹，虽说是姐妹们，到
        底是男女之分，日夜一处，起坐不方便，由不得叫人悬心。 既蒙老太太和太太的恩 典，把我派在二爷屋里，如今跟在园中住，都是我的干系。 太太想：多有无心中做
        出，有心人看见，当做有心事，反说坏了的，倒不如预先防着点儿。 况且二爷素日 的性格，太太是知道的，他又偏好在我们队里闹。}
}
{
    \eA{"I'll come again to see you to-morrow," she said, "but take good care of
        yourself! I gave the medicines I brought just now to Hsi Jen; let her rub
        you with them at night and I feel sure you'll get all right." With these
        recommendations, she walked out of the door. Hsi Jen hastened to catch her
        up and escorted her beyond the court. "Miss," she remarked, "we've really
        put you to the trouble of coming.}
}
{
}

\Verse{22}
{
    \zA{倘或不防，前后错了一点半点， 不论真假，人多嘴杂――那起坏人的嘴，太太还不知道呢：心顺了，说的比菩萨还 好； 心不顺，就没有忌讳了。
        二爷将来倘或有人说好，不过大家落个直过儿； 设若 叫人哼出一声不是来，我们不用说，粉身碎骨，还是平常，后来二爷一生的声名品 行，岂不完了呢？
        那时老爷太太也白疼了，白操了心了。 不如这会子防避些，似乎 妥当。 太太事情又多，一时固然想不到； 我们想不到便罢了，既想到了，要不回明
        了太太，罪越重了。}
}
{
    \eA{Some other day, when Mr. Secundus is well, I shall come in person to thank
        you." "What's there to thank me for?" replied Pao-ch'ai, turning her head
        round and smiling. "But mind, you advise him to carefully tend his health,
        and not to give way to idle thoughts and reckless ideas, and he'll recover.
        If there's anything he fancies to eat or to amuse himself with, come quietly
        over to me and fetch it for him.}
}
{
}

\Verse{23}
{
    \zA{近来我为这件事，日夜悬心，又恐怕太太听着生气，所以总没 敢言语。” 王夫人听了这话，正触了金钏儿之事，直呆了半晌，思前想后，心下越发感爱 袭人。
        笑道：“我的儿!你竟有这个心胸，想得这样周全。 我何曾又不想到这里？ 只 是这几次有事就混忘了。 你今日这话提醒了我，难为你这样细心，真真好孩子!也
        罢了，你且去罢，我自有道理。 只是还有一句话，你如今既说了这样的话，我索性 就把他交给你了。 好歹留点心儿，别叫他遭塌了身子才好。}
}
{
    \eA{There will be no use to disturb either our old lady, or Madame Wang, or any
        of the others; for in the event of its reaching Mr. Chia Cheng's ear,
        nothing may, at the time, come of it; but if by and bye he finds it to be
        true, we'll, doubtless, suffer for it!" While tendering this advice, she
        went on her way. Hsi Jen retraced her steps and returned into the room,
        fostering genuine feelings of gratitude for Pao-ch'ai.}
}
{
}

\Verse{24}
{
    \zA{自然不辜负你。” 袭人 低了一回头，方道：“太太吩咐，敢不尽心吗？” 说着，慢慢的退出。 回到院中，宝玉方醒。
        袭人回明香露之事，宝玉甚喜，即命调来吃，果然香妙 非常。 因心下惦着黛玉，要打发人去，只是怕袭人拦阻，便设法先使袭人往宝钗那 里去借书。
        袭人去了，宝玉便命晴雯来，吩咐道：“你到林姑娘那里，看他做什么 呢。 他要问我，只说我好了。” 晴雯道：“白眉赤眼儿的，作什么去呢!到底说句话
        儿，也像件事啊。”}
}
{
    \eA{But on entering, she espied Pao-yü silently lost in deep thought, and
        looking as if he were asleep, and yet not quite asleep, so she withdrew into
        the outer quarters to comb her hair and wash. Pao-yü meanwhile lay
        motionless in bed. His buttocks tingled with pain, as if they were pricked
        with needles, or dug with knives; giving him to boot a fiery sensation just
        as if fire were eating into them.}
}
{
}

\Verse{25}
{
    \zA{宝玉道：“没有什么可说的么。” 晴雯道：“或是送件东西，或是 取件东西，不然我去了怎么搭讪呢？” 宝玉想了一想，便伸手拿了两条旧绢子，撂
        与晴雯，笑道：“也罢，就说我叫你送这个给他去了。” 晴雯道：“这又奇了，他要 这半新不旧的两条绢子？ 他又要恼了，说你打趣他。”
        宝玉笑道：“你放心，他自然 知道。” 晴雯听了，只得拿了绢子，往潇湘馆来。 只见春纤正在栏杆上晾手巾，见他进 来，忙摇手儿说：“睡下了。”}
}
{
    \eA{He tried to change his position a bit, but unable to bear the anguish, he
        burst into groans. The shades of evening were by this time falling.
        Perceiving that though Hsi Jen had left his side there remained still two or
        three waiting-maids in attendance, he said to them, as he could find nothing
        for them to do just then, "You might as well go and comb your hair and
        perform your ablutions; come in, when I call you."}
}
{
}

\Verse{26}
{
    \zA{晴雯走进来，满屋漆黑，并未点灯，黛玉已睡在床上， 问：“是谁？” 晴雯忙答道：“晴雯。” 黛玉道：“做什么？” 晴雯道：“二爷叫给姑 娘送绢子来了。”
        黛玉听了，心中发闷，暗想：“做什么送绢子来给我？” 因问：“这 绢子是谁送他的？ 必定是好的，叫他留着送别人罢，我这会子不用这个。” 晴雯笑道：
        “不是新的，就是家常旧的。” 黛玉听了，越发闷住了。 细心揣度，一时方大悟过 来，连忙说：“放下，去罢。” 晴雯只得放下，抽身回去。}
}
{
    \eA{Hearing this, they likewise retired. During this while, Pao-yü fell into a
        drowsy state. Chiang Yü-han then rose before his vision and told him all
        about his capture by men from the Chung Shun mansion. Presently, Chin
        Ch'uan-erh too appeared in his room bathed in tears, and explained to him
        the circumstances which drove her to leap into the well.}
}
{
}

\Verse{27}
{
    \zA{一路盘算，不解何意。 这黛玉体贴出绢子的意思来，不觉神痴心醉，想到：宝玉能领会我这一番苦意， 又令我可喜。
        我这番苦意，不知将来可能如意不能，又令我可悲。 要不是这个意思， 忽然好好的送两块帕子来，竟又令我可笑了。 再想到私相传递，又觉可惧。 他既如
        此，我却每每烦恼伤心，反觉可愧。 如此左思右想，一时五内沸然。 由不得馀意缠 绵，便命掌灯，也想不起嫌疑避讳等事，研墨蘸笔，便向那两块旧帕上写道：
        眼空蓄泪泪空垂，暗洒闲抛更向谁？}
}
{
    \eA{But Pao-yü, who was half dreaming and half awake, was not able to give his
        mind to anything that was told him. Unawares, he became conscious of some
        one having given him a push; and faintly fell on his ear the plaintive tones
        of some person in distress. Pao-yü was startled out of his dreams. On
        opening his eyes, he found it to be no other than Lin Tai-yü.}
}
{
}

\Verse{28}
{
    \zA{尺幅鲛绡劳惠赠，为君那得不伤悲! 其　二 抛珠滚玉只偷潸，镇日无心镇日闲。 枕上袖边难拂拭，任他点点与斑斑。 其　三 彩线难收面上珠，湘江旧迹已模糊。
        窗前亦有千竿竹，不识香痕渍也无？ 那黛玉还要往下写时，觉得浑身火热，面上作烧，走至镜台揭起锦袱一照，只见腮 上通红，真合压倒桃花，却不知病由此起。
        一时方上床睡去，犹拿着绢子思索，不 在话下。 却说袭人来见宝钗，谁知宝钗不在园内，往他母亲那里去了。}
}
{
    \eA{But still fearing that it was only a dream, he promptly raised himself, and
        drawing near her face he passed her features under a minute scrutiny. Seeing
        her two eyes so swollen, as to look as big as peaches, and her face
        glistening all over with tears: "If it is not Tai-yü," (he thought), "who
        else can it be?"}
}
{
}

\Verse{29}
{
    \zA{袭人不便空手回 来，等至起更，宝钗方回。 原来宝钗素知薛蟠情性，心中已有一半疑是薛蟠挑唆了人来告宝玉了，谁知又 听袭人说出来，越发信了。
        究竟袭人是焙茗说的，那焙茗也是私心窥度，并未据实， 大家都是一半猜度，竟认作十分真切了。 可笑那薛蟠因素日有这个名声，其实这一
        次却不是他干的，竟被人生生的把个罪名坐定。 这日正从外头吃了酒回来，见过了 母亲，只见宝钗在这里坐着，说了几句闲话儿，忽然想起，因问道：“听见宝玉挨
        打，是为什么？”}
}
{
    \eA{Pao-yü meant to continue his scrutiny, but the lower part of his person gave
        him such unbearable sharp twitches that finding it a hard task to keep up,
        he, with a shout of "Ai-yo," lay himself down again, as he heaved a sigh.
        "What do you once more come here for?" he asked. "The sun, it is true, has
        set; but the heat remaining on the ground hasn't yet gone, so you may, by
        coming over, get another sunstroke.}
}
{
}

\Verse{30}
{
    \zA{薛姨妈正为这个不自在，见他问时，便咬着牙道：“不知好歹的 冤家，都是你闹的，你还有脸来问！” 薛蟠见说便怔了，忙问道：“我闹什么？” 薛
        姨妈道：“你还装腔呢!人人都知道是你说的。” 薛蟠道：“人人说我杀了人，也就信 了罢？” 薛姨妈道：“连你妹妹都知道是你说，难道他也赖你不成？”
        宝钗忙劝道： “妈妈和哥哥且别叫喊，消消停停的，就有个青红皂白了。” 又向薛蟠道：“是你说
        的也罢，不是你说的也罢，事情也过去了，不必较正，把小事倒弄大了。}
}
{
    \eA{Of course, I've had a thrashing but I don't feel any pains or aches. If I
        behave in this fashion, it's all put on to work upon their credulity, so
        that they may go and spread the reports outside in such a way as to reach my
        father's ear. Really it's all sham; so you mustn't treat it as a fact!"}
}
{
}

\Verse{31}
{
    \zA{我只劝你 从此以后少在外头胡闹，少管别人的事。 天天一处大家胡逛，你是个不防头的人， 过后没事就罢了，倘或有事，不是你干的，人人都也疑惑说是你干的。
        不用别人， 我先就疑惑你。” 薛蟠本是个心直口快的人，见不得这样藏头露尾的事； 又是宝钗劝他别再胡逛 去；
        他母亲又说他犯舌，宝玉之打，是他治的：早已急得乱跳，赌神发誓的分辩。 又骂众人：“谁这么编派我？ 我把那囚攮的牙敲了!分明是为打了宝玉，没的献勤儿，
        拿我来做幌子。 难道宝玉是天王？}
}
{
    \eA{Though Lin Tai-yü was not giving way at the time to any wails or loud sobs,
        yet the more she indulged in those suppressed plaints of hers, the worse she
        felt her breath get choked and her throat obstructed; so that when Pao-yü's
        assurances fell on her ear, she could not express a single sentiment, though
        she treasured thousands in her mind.}
}
{
}

\Verse{32}
{
    \zA{他父亲打他一顿，一家子定要闹几天。 那一回为 他不好，姨父打了他两下子，过后儿老太太不知怎么知道了，说是珍大哥治的，好 好儿的叫了去骂了一顿。
        今日越发拉上我了!既拉上我也不怕，索性进去把宝玉打 死了，我替他偿命！” 一面嚷，一面找起一根门闩来就跑。 慌的薛姨妈拉住骂道：“作
        死的孽障，你打谁去？ 你先打我来！” 薛蟠的眼急的铜铃一般，嚷道：“何苦来!又不 叫我去，为什么好好的赖我？}
}
{
    \eA{It was only after a long pause that she at last could observe, with agitated
        voice: "You must after this turn over a new leaf." At these words, Pao-yü
        heaved a deep sigh. "Compose your mind," he urged. "Don't speak to me like
        this; for I am quite prepared to even lay down my life for all those
        persons!"}
}
{
}

\Verse{33}
{
    \zA{将来宝玉活一日，我耽一日的口舌，不如大家死了清 净！” 宝钗忙也上前劝道：“你忍耐些儿罢。 妈妈急的这个样儿，你不说来劝，你倒 反闹的这样。
        别说是妈妈，就是旁人来劝你，也是为好，倒把你的性子劝上来！” 薛蟠道：“你这会子又说这话，都是你说的。” 宝钗道：“你只怨我说，再不怨你那
        顾前不顾后的形景！” 薛蟠道：“你只会怨我顾前不顾后，你怎么不怨宝玉外头招风 惹草的呢？}
}
{
    \eA{But scarcely had he concluded this remark than some one outside the court
        was heard to say: "Our lady Secunda has arrived." Lin Tai-yü readily
        concluded that it was lady Feng coming, so springing to her feet at once,
        "I'm off," she said; "out by the back-court. I'll look you up again by and
        bye." "This is indeed strange!" exclaimed Pao-yü as he laid hold of her and
        tried to detain her.}
}
{
}

\Verse{34}
{
    \zA{别说别的，就拿前日琪官儿的事比给你们听：那琪官儿我们见了十来次， 他并没和我说一句亲热话，怎么前儿他见了，连姓名还不知道，就把汗巾子给他？
        难道这也是我说的不成？” 薛姨妈和宝钗急的说道：“还提这个!可不是为这个打他 呢。 可见是你说的了。”
        薛蟠道：“真真的气死人了!赖我说的我不恼，我只气一个 宝玉闹的这么天翻地覆的！” 宝钗道：“谁闹来着？ 你先持刀动杖的闹起来，倒说别 人闹。”}
}
{
    \eA{"How is it that you've deliberately started living in fear and trembling of
        her!" Lin Tai-yü grew impatient and stamped her feet. "Look at my eyes!" she
        added in an undertone. "Must those people amuse themselves again by poking
        fun at me?" After this response, Pao-yü speedily let her go. Lin Tai-yü with
        hurried step withdrew behind the bed;}
}
{
}

\Verse{35}
{
    \zA{薛蟠见宝钗说的话句句有理，难以驳正，比母亲的话反难回答，因此便要设法 拿话堵回他去，就无人敢拦自己的话了。 也因正在气头儿上，未曾想话之轻重，便
        道：“好妹妹，你不用和我闹，我早知道你的心了。 从先妈妈和我说：你这金锁要 拣有玉的才可配，你留了心，见宝玉有那劳什子，你自然如今行动护着他。” 话未
        说了，把个宝钗气怔了，拉着薛姨妈哭道：“妈妈，你听哥哥说的是什么话！” 薛蟠 见妹子哭了，便知自己冒撞，便赌气走到自己屋里安歇不提。}
}
{
    \eA{and no sooner had she issued into the back-court, than lady Feng made her
        appearance in the room by the front entrance. "Are you better?" she asked
        Pao-yü. "If you fancy anything to eat, mind you send some one over to my
        place to fetch it for you." Thereupon Mrs. Hsüeh also came to pay him a
        visit. Shortly after, a messenger likewise arrived from old lady Chia (to
        inquire after him).}
}
{
}

\Verse{36}
{
    \zA{宝钗满心委屈气忿，待要怎样，又怕他母亲不安，少不得含泪别了母亲，各自 回来。 到屋里整哭了一夜。 次日一早起来，也无心梳洗，胡乱整理了衣裳，便出来
        瞧母亲。 可巧遇见黛玉独立在花阴之下，问他那里去，宝钗因说：“家去。” 口里说 着，便只管走。
        黛玉见他无精打彩的去了，又见眼上好似有哭泣之状，大非往日可 比，便在后面笑道：“姐姐也自己保重些儿。 就是哭出两缸泪来，也医不好棒疮！”
        不知宝钗如何答对，且听下回分解。 (本章完)}
}
{
    \eA{When the time came to prepare the lights, Pao-yü had a couple of mouthfuls
        of soup to eat, but he felt so drowsy and heavy that he fell asleep.
        Presently, Chou Jui's wife, Wu Hsin-teng's wife and Cheng Hao-shih's wife,
        all of whom were old dames who frequently went to and fro, heard that Pao-yü
        had been flogged and they too hurried into his quarters. Hsi Jen promptly
        went out to greet them.}
}
{
}

\Verse{37}
{
    \zA{}
}
{
    \eA{"Aunts," she whispered, smiling, "you've come a little too late; Master
        Secundus is sleeping." Saying this, she led them into the room on the
        opposite side, and, pressing then to sit down, she poured them some tea.
        After sitting perfectly still for a time, "When Master Secundus awakes" the
        dames observed, "do send us word!" Hsi Jen assured them that she would, and
        escorted them out. Just, however, as she was about to retrace her footsteps,
        she met an old matron, sent over by Madame Wang, who said to her: "Our
        mistress wants one of Master Secundus attendants to go and see her." Upon
        hearing this message, Hsi Jen communed with her own thoughts.}
}
{
}

\Verse{38}
{
    \zA{}
}
{
    \eA{Then turning round, she whispered to Ch'ing Wen, She Yüeh, Ch'iu Wen, and
        the other maids: "Our lady wishes to see one of us, so be careful and remain
        in the room while I go. I'll be back soon." At the close of her injunctions,
        she and the matron made their exit out of the garden by a short cut, and
        repaired into the drawing-room. Madame Wang was seated on the cool couch,
        waving a banana-leaf fan. When she became conscious of her arrival: "It
        didn't matter whom you sent," she remarked, "any one would have done. But
        have you left him again? Who's there to wait on him?" At this question, Hsi
        Jen lost no time in forcing a smile. "Master Secundus," she replied, "just
        now fell into a sound sleep.}
}
{
}

\Verse{39}
{
    \zA{}
}
{
    \eA{Those four or five girls are all right now, they are well able to attend to
        their master, so please, Madame, dispel all anxious thoughts! I was afraid
        that your ladyship might have some orders to give, and that if I sent any of
        them, they might probably not hear distinctly, and thus occasion delay in
        what there was to be done." "There's nothing much to tell you," added Madame
        Wang. "I only wish to ask how his pains and aches are getting on now?" "I
        applied on Mr. Secundus," answered Hsi Jen, "the medicine, which Miss
        Pao-ch'ai brought over; and he's better than he was. He was so sore at one
        time that he couldn't lie comfortably; but the deep sleep, in which he is
        plunged now, is a clear sign of his having improved." "Has he had anything
        to eat?" further inquired Madame Wang.}
}
{
}

\Verse{40}
{
    \zA{}
}
{
    \eA{"Our dowager mistress sent him a bowl of soup," Hsi Jen continued, "and of
        this he has had a few mouthfuls. He shouted and shouted that his mouth was
        parched and fancied a decoction of sour plums, but remembering that sour
        plums are astringent things, that he had been thrashed only a short time
        before, and that not having been allowed to groan, he must, of course, have
        been so hard pressed that fiery virus and heated blood must unavoidably have
        accumulated in the heart, and that were he to put anything of the kind
        within his lips, it might be driven into the cardiac regions and give rise
        to some serious illness;}
}
{
}

\Verse{41}
{
    \zA{}
}
{
    \eA{and what then would we do? I therefore reasoned with him for ever so long
        and at last succeeded in deterring him from touching any. So simply taking
        that syrup of roses, prepared with sugar, I mixed some with water and he had
        half a small cup of it. But he drank it with distaste; for, being surfeited
        with it, he found it neither scented nor sweet." "Ai-yah!" ejaculated Madame
        Wang. "Why didn't you come earlier and tell me? Some one sent me the other
        day several bottles of scented water. I meant at one time to have given him
        some, but as I feared that it would be mere waste, I didn't let him have
        any. But since he is so sick and tired of that preparation of roses, that he
        turns up his nose at it, take those two bottles with you.}
}
{
}

\Verse{42}
{
    \zA{}
}
{
    \eA{If you just mix a teaspoonful of it in a cup of water, it will impart to it
        a very strong perfume." So saying, she hastened to tell Ts'ai Yün to fetch
        the bottles of scented water, which she had received as a present a few days
        before. "Let her only bring a couple of them, they'll be enough!" Hsi Jen
        chimed in. "If you give us more, it will be a useless waste! If it isn't
        enough, I can come and fetch a fresh supply. It will come to the same
        thing!" Having listened to all they had to say, Ts'ai Yün left the room.
        After some considerable time, she, in point of fact, returned with only a
        couple of bottles, which she delivered to Hsi Jen. On examination, Hsi Jen
        saw two small glass bottles, no more than three inches in size, with
        screwing silver stoppers at the top.}
}
{
}

\Verse{43}
{
    \zA{}
}
{
    \eA{On the gosling-yellow labels was written, on one: "Pure extract of \_olea
        fragrans\_," on the other, "Pure extract of roses." "What fine things these
        are!" Hsi Jen smiled. "How many small bottles the like of this can there
        be?" "They are of the kind sent to the palace," rejoined Madame Wang.
        "Didn't you notice that gosling-yellow slip? But mind, take good care of
        them for him; don't fritter them away!" Hsi Jen assented. She was about to
        depart when Madame Wang called her back. "I've thought of something," she
        said, "that I want to ask you." Hsi Jen hastily came back. Madame Wang made
        sure that there was no one in the room. "I've heard a faint rumour," she
        then inquired, "to the effect that Pao-yü got a thrashing on this occasion
        on account of something or other which Huan-Erh told my husband.}
}
{
}

\Verse{44}
{
    \zA{}
}
{
    \eA{Have you perchance heard what it was that he said? If you happen to learn
        anything about it, do confide in me, and I won't make any fuss and let
        people know that it was you who told me." "I haven't heard anything of the
        kind," answered Hsi Jen. "It was because Mr. Secundus forcibly detained an
        actor, and that people came and asked master to restore him to them that he
        got flogged." "It was also for this," continued Madame Wang as she nodded
        her head, "but there's another reason besides." "As for the other reason, I
        honestly haven't the least idea about it," explained Hsi Jen. "But I'll make
        bold to-day, and say something in your presence, Madame, about which I don't
        know whether I am right or wrong in speaking. According to what's proper…."}
}
{
}

\Verse{45}
{
    \zA{}
}
{
    \eA{She had only spoken half a sentence, when hastily she closed her mouth
        again. "You are at liberty to proceed," urged Madame Wang. "If your ladyship
        will not get angry, I'll speak out," remarked Hsi Jen. "Why should I get
        angry?" observed Madame Wang. "Proceed!" "According to what's proper,"
        resumed Hsi Jen, "our Mr. Secundus should receive our master's admonition,
        for if master doesn't hold him in check, there's no saying what he mightn't
        do in the future." As soon as Madame Wang heard this, she clasped her hands
        and uttered the invocation, "O-mi-to-fu!" Unable to resist the impulse, she
        drew near Hsi Jen.}
}
{
}

\Verse{46}
{
    \zA{}
}
{
    \eA{"My dear child," she added, "you have also luckily understood the real state
        of things. What you told me is in perfect harmony with my own views! Is it
        likely that I don't know how to look after a son? In former days, when your
        elder master, Chu, was alive, how did I succeed in keeping him in order? And
        can it be that I don't, after all, now understand how to manage a son? But
        there's a why and a wherefore in it. The thought is ever present in my mind
        now, that I'm already a woman past fifty, that of my children there only
        remains this single one, that he too is developing a delicate physique, and
        that, what's more, our dear senior prizes him as much as she would a jewel,
        that were he kept under strict control, and anything perchance to happen to
        him, she might, an old lady as she is, sustain some harm from resentment,
        and that as the high as well as the low will then have no peace or quiet,
        won't things get in a bad way?}
}
{
}

\Verse{47}
{
    \zA{}
}
{
    \eA{So I feel prompted to spoil him by over-indulgence. Time and again I reason
        with him. Sometimes, I talk to him; sometimes, I advise him; sometimes, I
        cry with him. But though, for the time being, he's all right, he doesn't,
        later on, worry his mind in any way about what I say, until he positively
        gets into some other mess, when he settles down again. But should any harm
        befall him, through these floggings, upon whom will I depend by and bye?" As
        she spoke, she could not help melting into tears. At the sight of Madame
        Wang in this disconsolate mood, Hsi Jen herself unconsciously grew wounded
        at heart, and as she wept along with her, "Mr. Secundus," she ventured, "is
        your ladyship's own child, so how could you not love him?}
}
{
}

\Verse{48}
{
    \zA{}
}
{
    \eA{Even we, who are mere servants, think it a piece of good fortune when we can
        wait on him for a time, and all parties can enjoy peace and quiet. But if he
        begins to behave in this manner, even peace and quiet will be completely out
        of the question for us. On what day, and at what hour, don't I advise Mr.
        Secundus; yet I can't manage to stir him up by any advice! But it happens
        that all that crew are ever ready to court his friendship, so it isn't to be
        wondered that he is what he is! The truth is that he thinks the advice we
        give him is not right and proper! As you have to-day, Madame, alluded to
        this subject, I've got something to tell you which has weighed heavy on my
        mind.}
}
{
}

\Verse{49}
{
    \zA{}
}
{
    \eA{I've been anxious to come and confide it to your ladyship and to solicit
        your guidance, but I've been in fear and dread lest you should give way to
        suspicion. For not only would then all my disclosures have been in vain, but
        I would have deprived myself of even a piece of ground wherein my remains
        could be laid." Madame Wang perceived that her remarks were prompted by some
        purpose. "My dear child," she eagerly urged; "go on, speak out! When I
        recently heard one and all praise you secretly behind your back, I simply
        fancied that it was because you were careful in your attendance on Pao-yü;
        or possibly because you got on well with every one; all on account of minor
        considerations like these; (but I never thought it was on account of your
        good qualities).}
}
{
}

\Verse{50}
{
    \zA{}
}
{
    \eA{As it happens, what you told me just now concerns, in all its bearings, a
        great principle, and is in perfect accord with my ideas, so speak out
        freely, if you have aught to say! Only let no one else know anything about
        it, that is all that is needed." "I've got nothing more to say," proceeded
        Hsi Jen. "My sole idea was to solicit your advice, Madame, as to how to
        devise a plan to induce Mr. Secundus to move his quarters out of the garden
        by and bye, as things will get all right then." This allusion much alarmed
        Madame Wang. Speedily taking Hsi Jen's hand in hers: "Is it likely," she
        inquired, "that Pao-yü has been up to any mischief with any one?" "Don't be
        too suspicious!" precipitately replied Hsi Jen.}
}
{
}

\Verse{51}
{
    \zA{}
}
{
    \eA{"It wasn't at anything of the kind that I was hinting. I merely expressed my
        humble opinion. Mr. Secundus is a young man now, and the young ladies inside
        are no more children. More than that, Miss Lin and Miss Pao may be two
        female maternal first cousins of his, but albeit his cousins, there is
        nevertheless the distinction of male and female between them; and day and
        night, as they are together, it isn't always convenient, when they have to
        rise and when they have to sit; so this cannot help making one give way to
        misgivings. Were, in fact, any outsider to see what's going on, it would not
        look like the propriety, which should exist in great families. The proverb
        appositely says that: 'when there's no trouble, one should make provision
        for the time of trouble.'}
}
{
}

\Verse{52}
{
    \zA{}
}
{
    \eA{How many concerns there are in the world, of which there's no making head or
        tail, mostly because what persons do without any design is construed by such
        designing people, as chance to have their notice attracted to it, as having
        been designedly accomplished, and go on talking and talking till, instead of
        mending matters, they make them worse! But if precautions be not taken
        beforehand, something improper will surely happen, for your ladyship is well
        aware of the temperament Mr. Secundus has shown all along!}
}
{
}

\Verse{53}
{
    \zA{}
}
{
    \eA{Besides, his great weakness is to fuss in our midst, so if no caution be
        exercised, and the slightest mistake be sooner or later committed, there'll
        be then no question of true or false: for when people are many one says one
        thing and another, and what is there that the mouths of that mean lot will
        shun with any sign of respect? Why, if their hearts be well disposed, they
        will maintain that he is far superior to Buddha himself. But if their hearts
        be badly disposed, they will at once knit a tissue of lies to show that he
        cannot even reach the standard of a beast! Now, if people by and bye speak
        well of Mr. Secundus, we'll all go on smoothly with our lives.}
}
{
}

\Verse{54}
{
    \zA{}
}
{
    \eA{But should he perchance give reason to any one to breathe the slightest
        disparaging remark, won't his body, needless for us to say, be smashed to
        pieces, his bones ground to powder, and the blame, which he might incur, be
        made ten thousand times more serious than it is? These things are all
        commonplace trifles; but won't Mr. Secundus' name and reputation be
        subsequently done for for life? Secondly, it's no easy thing for your
        ladyship to see anything of our master. A proverb also says: 'The perfect
        man makes provision beforehand;' so wouldn't it be better that we should,
        this very minute, adopt such steps as will enable us to guard against such
        things? Your ladyship has much to attend to, and you couldn't, of course,
        think of these things in a moment.}
}
{
}

\Verse{55}
{
    \zA{}
}
{
    \eA{And as for us, it would have been well and good, had they never suggested
        themselves to our minds; but since they have, we should be the more to blame
        did we not tell you anything about them, Madame. Of late, I have racked my
        mind, both day and night on this score; and though I couldn't very well
        confide to any one, my lamp alone knows everything!" After listening to
        these words, Madame Wang felt as if she had been blasted by thunder and
        struck by lightning; and, as they fitted so appositely with the incident
        connected with Chin Ch'uan-erh, her heart was more than ever fired with
        boundless affection for Hsi Jen. "My dear girl," she promptly smiled, "it's
        you, who are gifted with enough foresight to be able to think of these
        things so thoroughly.}
}
{
}

\Verse{56}
{
    \zA{}
}
{
    \eA{Yet, did I not also think of them? But so busy have I been these several
        times that they slipped from my memory. What you've told me to-day, however,
        has brought me to my senses! It's, thanks to you, that the reputation of me,
        his mother, and of him, my son, is preserved intact! I really never had the
        faintest idea that you were so excellent! But you had better go now; I know
        of a way. Yet, just another word. After your remarks to me, I'll hand him
        over to your charge; please be careful of him. If you preserve him from
        harm, it will be tantamount to preserving me from harm, and I shall
        certainly not be ungrateful to you for it." Hsi Jen said several consecutive
        yes's, and went on her way. She got back just in time to see Pao-yü awake.
        Hsi Jen explained all about the scented water;}
}
{
}

\Verse{57}
{
    \zA{}
}
{
    \eA{and, so intensely delighted was Pao-yü, that he at once asked that some
        should be mixed and brought to him to taste. In very deed, he found it
        unusually fragrant and good. But as his heart was a prey to anxiety on
        Tai-yü's behalf, he was full of longings to despatch some one to look her
        up. He was, however, afraid of Hsi Jen. Readily therefore he devised a plan
        to first get Hsi Jen out of the way, by despatching her to Pao-ch'ai's, to
        borrow a book. After Hsi Jen's departure, he forthwith called Ch'ing Wen.
        "Go," he said, "over to Miss Lin's and see what she's up to. Should she
        inquire about me, all you need tell her is that I'm all right." "What shall
        I go empty-handed for?" rejoined Ch'ing Wen.}
}
{
}

\Verse{58}
{
    \zA{}
}
{
    \eA{"If I were, at least, to give her a message, it would look as if I had gone
        for something." "I have no message that you can give her," added Pao-yü. "If
        it can't be that," suggested Ch'ing Wen; "I might either take something over
        or fetch something. Otherwise, when I get there, what excuse will I be able
        to find?" After some cogitation, Pao-yü stretched out his hand and, laying
        hold of a couple of handkerchiefs, he threw them to Ch'ing Wen. "These will
        do," he smiled. "Just tell her that I bade you take them to her." "This is
        strange!" exclaimed Ch'ing Wen. "Will she accept these two half worn-out
        handkerchiefs! She'll besides get angry and say that you were making fun of
        her." "Don't worry yourself about that;"}
}
{
}

\Verse{59}
{
    \zA{}
}
{
    \eA{laughed Pao-yü. "She will certainly know what I mean." Ch'ing Wen, at this
        rejoinder, had no help but to take the handkerchiefs and to go to the Hsiao
        Hsiang lodge, where she discovered Ch'un Hsien in the act of hanging out
        handkerchiefs on the railings to dry. As soon as she saw her walk in, she
        vehemently waved her hand. "She's gone to sleep!" she said. Ch'ing Wen,
        however, entered the room. It was in perfect darkness. There was not even so
        much as a lantern burning, and Tai-yü was already ensconced in bed. "Who is
        there?" she shouted. "It's Ch'ing Wen!" promptly replied Ch'ing Wen. "What
        are you up to?" Tai-yü inquired. "Mr. Secundus," explained Ch'ing Wen,
        "sends you some handkerchiefs, Miss." Tai-yü's spirits sunk as soon as she
        caught her reply. "What can he have sent me handkerchiefs for?"}
}
{
}

\Verse{60}
{
    \zA{}
}
{
    \eA{she secretly reasoned within herself. "Who gave him these handkerchiefs?"
        she then asked aloud. "They must be fine ones, so tell him to keep them and
        give them to some one else; for I don't need such things at present."
        "They're not new," smiled Ch'ing Wen. "They are of an ordinary kind, and
        old." Hearing this, Lin Tai-yü felt downcast. But after minutely searching
        her heart, she at last suddenly grasped his meaning and she hastily
        observed: "Leave them and go your way." Ch'ing Wen was compelled to put them
        down; and turning round, she betook herself back again. But much though she
        turned things over in her mind during the whole of her way homewards, she
        did not succeed in solving their import.}
}
{
}

\Verse{61}
{
    \zA{}
}
{
    \eA{When Tai-yü guessed the object of the handkerchief, her very soul unawares
        flitted from her. "As Pao-yü has gone to such pains," she pondered, "to try
        and probe this dejection of mine, I have, on one hand, sufficient cause to
        feel gratified; but as there's no knowing what my dejection will come to in
        the future there is, on the other, enough to make me sad. Here he abruptly
        and deliberately sends me a couple of handkerchiefs; and, were it not that
        he has divined my inmost feelings, the mere sight of these handkerchiefs
        would be enough to make me treat the whole thing as ridiculous. The secret
        exchange of presents between us," she went on to muse, "fills me also with
        fears;}
}
{
}

\Verse{62}
{
    \zA{}
}
{
    \eA{and the thought that those tears, which I am ever so fond of shedding to
        myself, are of no avail, drives me likewise to blush with shame." And by
        dint of musing and reflecting, her heart began, in a moment, to bubble over
        with such excitement that, much against her will, her thoughts in their
        superabundance rolled on incessantly. So speedily directing that a lamp
        should be lighted, she little concerned herself about avoiding suspicion,
        shunning the use of names, or any other such things, and set to work and
        rubbed the ink, soaked the pen, and then wrote the following stanzas on the
        two old handkerchiefs:  Vain in my eyes the tears collect; those tears in
        vain they flow,  Which I in secret shed; they slowly drop; but for whom
        though?}
}
{
}

\Verse{63}
{
    \zA{}
}
{
    \eA{The silk kerchiefs, which he so kindly troubled to give me,  How ever could
        they not with anguish and distress fill me? The second ran thus:  Like
        falling pearls or rolling gems, they trickle on the sly. Daily I have no
        heart for aught; listless all day am I.  As on my pillow or sleeves' edge I
        may not wipe them dry,  I let them dot by dot, and drop by drop to run
        freely. And the third:  The coloured thread cannot contain the pearls
        cov'ring my face. Tears were of old at Hsiang Chiang shed, but faint has
        waxed each  trace. Outside my window thousands of bamboos, lo, also grow,
        But whether they be stained with tears or not, I do not know.}
}
{
}

\Verse{64}
{
    \zA{}
}
{
    \eA{Lin Tai-yü was still bent upon going on writing, but feeling her whole body
        burn like fire, and her face scalding hot, she advanced towards the
        cheval-glass, and, raising the embroidered cover, she looked in. She saw at
        a glance that her cheeks wore so red that they, in very truth, put even the
        peach blossom to the shade. Yet little did she dream that from this date her
        illness would assume a more serious phase. Shortly, she threw herself on the
        bed, and, with the handkerchiefs still grasped in her hand, she was lost in
        a reverie. Putting her aside, we will now take up our story with Hsi Jen.
        She went to pay a visit to Pao-ch'ai, but as it happened, Pao-ch'ai was not
        in the garden, but had gone to look up her mother.}
}
{
}

\Verse{65}
{
    \zA{}
}
{
    \eA{Hsi Jen, however, could not very well come back with empty hands so she
        waited until the second watch, when Pao-ch'ai eventually returned to her
        quarters. Indeed, so correct an estimate of Hsüeh P'an's natural disposition
        did Pao-ch'ai ever have, that from an early moment she entertained within
        herself some faint suspicion that it must have been Hsüeh P'an, who had
        instigated some person or other to come and lodge a complaint against
        Pao-yü. And when she also unexpectedly heard Hsi Jen's disclosures on the
        subject, she became more positive in her surmises.}
}
{
}

\Verse{66}
{
    \zA{}
}
{
    \eA{The one, who had, in fact, told Hsi Jen was Pei Ming. But Pei Ming too had
        arrived at the conjecture in his own mind, and could not adduce any definite
        proof, so that every one treated his statements as founded partly on mere
        suppositions, and partly on actual facts; but, despite this, they felt quite
        certain that it was (Hsüeh P'an) who had intrigued. Hsüeh P'an had always
        enjoyed this reputation; but on this particular instance the harm was not,
        actually, his own doing; yet as every one, with one consent, tenaciously
        affirmed that it was he, it was no easy matter for him, much though he might
        argue, to clear himself of blame. Soon after his return, on this day, from a
        drinking bout out of doors, he came to see his mother; but finding Pao-ch'ai
        in her rooms, they exchanged a few irrelevant remarks.}
}
{
}

\Verse{67}
{
    \zA{}
}
{
    \eA{"I hear," he consequently asked, "that cousin Pao-yü has got into trouble;
        why is it?" Mrs. Hsüeh was at the time much distressed on this score. As
        soon therefore as she caught this question, she gnashed her teeth with rage,
        and shouted: "You good-for-nothing spiteful fellow! It's all you who are at
        the bottom of this trouble; and do you still have the face to come and ply
        me with questions?" These words made Hsüeh P'an wince. "When did I stir up
        any trouble?" he quickly asked. "Do you still go on shamming!" cried Mrs.
        Hsüeh. "Every one knows full well that it was you, who said those things,
        and do you yet prevaricate?" "Were every one," insinuated Hsüeh P'an, "to
        assert that I had committed murder, would you believe even that?" "Your very
        sister is well aware that they were said by you."}
}
{
}

\Verse{68}
{
    \zA{}
}
{
    \eA{Mrs. Hsüeh continued, "and is it likely that she would accuse you falsely,
        pray?" "Mother," promptly interposed Pao-ch'ai, "you shouldn't be brawling
        with brother just now! If you wait quietly, we'll find out the plain and
        honest truth." Then turning towards Hsüeh P'an: "Whether it's you, who said
        those things or not," she added, "it's of no consequence. The whole affair,
        besides, is a matter of the past, so what need is there for any arguments;
        they will only be making a mountain of a mole-hill! I have just one word of
        advice to give you; don't, from henceforward, be up to so much reckless
        mischief outside; and concern yourself a little less with other people's
        affairs! All you do is day after day to associate with your friends and
        foolishly gad about!}
}
{
}

\Verse{69}
{
    \zA{}
}
{
    \eA{You are a happy-go-lucky sort of creature! If nothing happens well and good;
        but should by and bye anything turn up, every one will, though it be none of
        your doing, imagine again that you are at the bottom of it! Not to speak of
        others, why I myself will be the first to suspect you!" Hsüeh P'an was
        naturally open-hearted and plain-spoken, and could not brook anything in the
        way of innuendoes, so, when on the one side, Pao-ch'ai advised him not to
        foolishly gad about, and his mother, on the other, hinted that he had a foul
        tongue, and that he was the cause that Pao-yü had been flogged, he at once
        got so exasperated that he jumped about in an erratic manner and did all in
        his power, by vowing and swearing, to explain matters.}
}
{
}

\Verse{70}
{
    \zA{}
}
{
    \eA{"Who has," he ejaculated, heaping abuse upon every one, "laid such a tissue
        of lies to my charge! I'd like to take the teeth of that felon and pull them
        out! It's clear as day that they shove me forward as a target; for now that
        Pao-yü has been flogged they find no means of making a display of their
        zeal. But, is Pao-yü forsooth the lord of the heavens that because he has
        had a thrashing from his father, the whole household should be fussing for
        days? The other time, he behaved improperly, and my uncle gave him two
        whacks.}
}
{
}

\Verse{71}
{
    \zA{}
}
{
    \eA{But our venerable ancestor came, after a time, somehow or other, I don't
        know how, to hear about it, and, maintaining that it was all due to Mr. Chia
        Chen, she called him before her, and gave him a good blowing up. And here
        to-day, they have gone further, and involved me. They may drag me in as much
        as they like, I don't fear a rap! But won't it be better for me to go into
        the garden, and take Pao-yü and give him a bit of my mind and kill him? I
        can then pay the penalty by laying down my life for his, and one and all
        will enjoy peace and quiet!" While he clamoured and shouted, he looked about
        him for the bar of the door, and, snatching it up, he there and then was
        running off, to the consternation of Mrs. Hsüeh, who clutched him in her
        arms.}
}
{
}

\Verse{72}
{
    \zA{}
}
{
    \eA{"You murderous child of retribution!" she cried. "Whom would you go and
        beat? come first and assail me?" From excitement Hsüeh P'an's eyes protruded
        like copper bells. "What are you up to," he vociferated, "that you won't let
        me go where I please, and that you deliberately go on calumniating me? But
        every day that Pao-yü lives, the longer by that day I have to bear a false
        charge, so it's as well that we should both die that things be cleared up?"
        Pao-ch'ai too hurriedly rushed forward. "Be patient a bit!" she exhorted
        him. "Here's mamma in an awful state of despair. Not to mention that it
        should be for you to come and pacify her, you contrariwise kick up all this
        rumpus!}
}
{
}

\Verse{73}
{
    \zA{}
}
{
    \eA{Why, saying nothing about her who is your parent, were even a perfect
        stranger to advise you, it would be meant for your good! But the good
        counsel she gave you has stirred up your monkey instead." "From the way
        you're now speaking," Hsüeh P'an rejoined, "it must be you, who said that it
        was I; no one else but you!" "You simply know how to feel displeased with me
        for speaking," argued Pao-ch'ai, "but you don't feel displeased with
        yourself for that reckless way of yours of looking ahead and not minding
        what is behind!" "You now bear me a grudge," Hsüeh P'an added, "for looking
        to what is ahead and not to what is behind; but how is it you don't feel
        indignant with Pao-yü for stirring up strife and provoking trouble outside?}
}
{
}

\Verse{74}
{
    \zA{}
}
{
    \eA{Leaving aside everything else, I'll merely take that affair of Ch'i
        Kuan-erh's, which occurred the other day, and recount it to you as an
        instance. My friends and I came across this Ch'i Kuan-erh, ten times at
        least, but never has he made a single intimate remark to me, and how is it
        that, as soon as he met Pao-yü the other day, he at once produced his sash,
        and gave it to him, though he did not so much as know what his surname and
        name were? Now is it likely, forsooth, that this too was something that I
        started?" "Do you still refer to this?"}
}
{
}

\Verse{75}
{
    \zA{}
}
{
    \eA{exclaimed Mrs. Hsüeh and Pao-ch'ai, out of patience. "Wasn't it about this
        that he was beaten? This makes it clear enough that it's you who gave the
        thing out." "Really, you're enough to exasperate one to death!" Hsüeh P'an
        exclaimed. "Had you confined yourselves to saying that I had started the
        yarn, I wouldn't have lost my temper; but what irritates me is that such a
        fuss should be made for a single Pao-yü, as to subvert heaven and earth!"
        "Who fusses?" shouted Pao-ch'ai. "You are the first to arm yourself to the
        teeth and start a row, and then you say that it's others who are up to
        mischief!"}
}
{
}

\Verse{76}
{
    \zA{}
}
{
    \eA{Hsüeh P'an, seeing that every remark, made by Pao-ch'ai, contained so much
        reasonableness that he could with difficulty refute it, and that her words
        were even harder for him to reply to than were those uttered by his mother,
        he was consequently bent upon contriving a plan to make use of such language
        as could silence her and compel her to return to her room, so as to have no
        one bold enough to interfere with his speaking; but, his temper being up, he
        was not in a position to weigh his speech. "Dear Sister!" he readily
        therefore said, "you needn't be flying into a huff with me! I've long ago
        divined your feelings. Mother told me some time back that for you with that
        gold trinket, must be selected some suitor provided with a jade one; as such
        a one will be a suitable match for you.}
}
{
}

\Verse{77}
{
    \zA{}
}
{
    \eA{And having treasured this in your mind, and seen that Pao-yü has that
        rubbishy thing of his, you naturally now seize every occasion to screen
        him…." However, before he could finish, Pao-ch'ai trembled with anger, and
        clinging to Mrs. Hsüeh, she melted into tears. "Mother," she observed, "have
        you heard what brother says, what is it all about?" Hsüeh P'an, at the sight
        of his sister bathed in tears, became alive to the fact that he had spoken
        inconsiderately, and, flying into a rage, he walked away to his own quarters
        and retired to rest. But we can well dispense with any further comment on
        the subject. Pao-ch'ai was, at heart, full of vexation and displeasure.}
}
{
}

\Verse{78}
{
    \zA{}
}
{
    \eA{She meant to give vent to her feelings in some way, but the fear again of
        upsetting her mother compelled her to conceal her tears. She therefore took
        leave of her parent, and went back all alone. On her return to her chamber,
        she sobbed and sobbed throughout the whole night. The next day, she got out
        of bed, as soon as it dawned; but feeling even no inclination to comb her
        chevelure or perform her ablutions, she carelessly adjusted her clothes and
        came out of the garden to see her mother. As luck would have it, she
        encountered Tai-yü standing alone under the shade of the trees, who inquired
        of her: "Where she was off to?" "I'm going home," Hsüeh Pao-ch'ai replied.
        And as she uttered these words, she kept on her way.}
}
{
}

\Verse{79}
{
    \zA{}
}
{
    \eA{But Tai-yü perceived that she was going off in a disconsolate mood; and,
        noticing that her eyes betrayed signs of crying, and that her manner was
        unlike that of other days, she smilingly called out to her from behind:
        "Sister, you should take care of yourself a bit. Were you even to cry so
        much as to fill two water jars with tears, you wouldn't heal the wounds
        inflicted by the cane." But as what reply Hsüeh Pao-ch'ai gave is not yet
        known to you, reader, lend an ear to the explanation contained in the next
        chapter.}
}
{
}
